%\chapter{Reference guide}

\section{Coupling with the 2D magnetic equilibrium}
In this section it is detailed the scheme which Astra--7.0 uses to couple with a
generic Grad--Shafranov solver, be it prescribed or free--boundary.
The main idea behind the coupling scheme is: the use of normalized grid $x$
together with the current diffusion being solved for $\psi$, suggest to devise a
scheme in which $\psi$ becomes the fundamental grid. This scheme is somewhat 
different from what is done in so--called $q$--solvers, which perform iterations
around the invariance of the safety factor $q$. Other schemes use the parallel 
current $j_{||}$ as the quantity that directs the iterations. 
The common experience is that both $q$--solvers and $j_{||}$--solvers are prone
to numerical instabilities, hampering the ease of simulating disparate
scenarios.

In the following it is explained why, what is called later on, $\psi$--solver is
better suited to guarantee stability and accuracy by resorting to not so 
complicated numerical scheme.
%
\subsection{The Grad--Shafranov equation for a $\psi$--solver} 
The starting point is the Grad--Shafranov equation (GSE):
\begin{eqnarray}
\Delta^* \Psi = -4\pi^2 \mu_0 R^2 \frac{\partial P}{\partial\psi} - 4\pi^2 F  \frac{\partial F}{\partial\psi} + ...{\rm (coils)}...
\label{GSeqcouplsec}
\end{eqnarray}
where $\Psi(R,Z)$, $\psi(x)$, and one requires that $\Psi(x(R,Z))=\psi(x)$ on closed flux surfaces.
$F=R B_\phi$ is the usual poloidal current integral.

A useful relation is found between $F$ and derivative of $\psi$ over $x$:
\begin{eqnarray}
F(x) = \frac{2\pi}{g_3} \Phi_{\rm b} H \hat{q} \nonumber \\
H = \frac{\partial\psi}{\partial V} \nonumber \\
\hat{q} = 2x\frac{\partial x}{\partial\psi} 
\label{Ffuncqhateq}
\end{eqnarray}
where $\displaystyle g_3=\left\langle \frac{1}{R^2}\right\rangle$. Notice that in the definition of $F$ 
there are two unknowns at the level of transport: $H(x)$, and $\Phi_{\rm b}$. These two quantities cannot 
be determined solely from 1D radial transport processes. However they are essential to compute both the 
poloidal (and thus the parallel) current, and to compute the radial coordinate $\rho=\rho_{\rm b} x$ that 
is actually used
in the transport equations. The user can also prescribe the value of $F$ at the plasma/vacuum interface, 
i.e. $F_{\rm b}=F_{\rm vac}$. If there are no poloidal currents flowing outside the separatrix, then 
$F_{\rm b}=R_0 B_{\phi.0}$, i.e. the vacuum constant value. The value of $F_{\rm vac}$ {\it per se} does 
not influence the equilibrium solution. However it influences the value of $\rho_{\rm b}$.

Note also that $P'=\partial P/\partial\psi$ is a known function of $x$.

Upon substitution of definition (\ref{Ffuncqhateq}) into the GSE we get (external coil currents are not 
mentioned from now on):
\begin{eqnarray}
\Delta^* \Psi = -4\pi^2 \mu_0 R^2 \frac{\partial P}{\partial\psi} - 16\pi^4 \Phi_{\rm b}^2\frac{\hat{q} }{g_3}  H  \frac{\partial }{\partial\psi} \left(\frac{\hat{q}}{g_3} H \right)
\label{GSeqcouplsec}
\end{eqnarray}

The flux--surface average of equation (\ref{GSeqcouplsec}) looks like:
\begin{eqnarray}
\frac{\partial}{\partial V}\left(g_2 H \right) = -4\pi^2 \mu_0 \frac{\partial P}{\partial\psi} - 16\pi^4 \Phi_{\rm b}^2 \hat{q}  H  \frac{\partial }{\partial\psi} \left(\frac{\hat{q}}{g_3} H \right)
\label{GSeqcouplsecfsa}
\end{eqnarray}
with $\displaystyle g_2=\left\langle \frac{|\nabla V|^2}{R^2}\right\rangle$.

Or:
\begin{eqnarray}
H\frac{\partial}{\partial x}\left(g_2 H \right) = -4\pi^2 \mu_0 \frac{\partial P}{\partial x} - 16\pi^4 \Phi_{\rm b}^2 \hat{q}  H  \frac{\partial }{\partial x} \left(\frac{\hat{q}}{g_3} H \right)
\label{GSeqcouplsecfsax}
\end{eqnarray}

By defining $P_x = \partial P / \partial x$, and the additional quantities: $y=H^2/2$, $A=g_2+16 \pi^4 \Phi_{\rm b}^2 \hat{q}^2/g_3$, $B=2(g_{2,x}+16 \pi^4 \Phi_{\rm b}^2 \hat{q}(\hat{q}/g_3)_x)$, and $C=-4\pi^2 \mu_0 P_x$, we can compact the equation in:
\begin{eqnarray}
A y_x = -B y + C
\label{GSeqcouplsecfsay}
\end{eqnarray}
which solution is easily computed as:
\begin{eqnarray}
y(x) =y_0 e^{-[B](x)} + e^{-[B](x)}\int_0^x{e^{[B](x')}C(x')dx'} 
\label{GSeqcouplsecfsaysol1}
\end{eqnarray}
where here we temporarily employ the notation $[B](x)=\int_0^x Bdx$.

The boundary condition $y(x=1)=y(1)$ is given in terms of known quantities:
\begin{eqnarray}
y(1) = \frac{1}{2}\left(\frac{F_{\rm vac} g_{3}(1)}{2\pi\Phi_{\rm b}\hat{q}(1)}\right)^2 
\label{y1bc}
\end{eqnarray}
which allows to compute the solution to equation (\ref{GSeqcouplsecfsaysol1}):
\begin{eqnarray}
y(x) =\frac{y(1)-e^{-[B](1)}\int_0^1{e^{[B](x')}C(x')dx'}}{e^{[B](1)}} e^{-[B](x)} + e^{-[B](x)}\int_0^x{e^{[B](x')}C(x')dx'} 
\label{GSeqcouplsecfsaysol2}
\end{eqnarray}
Having this solution, we can finally solve for $H(x)=\sqrt{2y(x)}$.

The knowledge of $H(x)$ allows to compute the plasma volume $V(x)$ as:
\begin{eqnarray}
V(x) =\int_0^x \frac{2x}{H\hat{q}}dx 
\label{VfromHeq}
\end{eqnarray}
On the other hand, the total plasma volume $V_{\rm b}=V(1)$ cannot be solved for by the flux--surface
averaged GSE, since it depends on the exact last closed flux surface shape. 

What then happens is that equation (\ref{VfromHeq}) is actually a non--linear equation for $\Phi_{\rm b}$:
\begin{eqnarray}
V_{\rm b}(\Phi_b) =\int_0^1 \frac{2x}{H(\Phi_{\rm b})\hat{q}}dx 
\label{VfromHeq}
\end{eqnarray}

Having now found both $H(x)$ and the $\Phi_{\rm b}$ consistent with $\psi(x)$ and with the plasma volume
$V_{\rm b}$, we can solve the 2D problem, equation (\ref{GSeqcouplsec}), with boundary condition
$\Psi_{\infty}=0$. The 2D problem is solved in this form now:
\begin{eqnarray}
\Delta^* \Psi = -4\pi^2 \mu_0 \left(R^2-\frac{1}{g_3}\right) \frac{\partial P}{\partial\psi} + \frac{1 }{g_3}  \frac{\partial }{\partial V} \left(g_2 H \right)
\label{GSeqcouplsecfinal}
\end{eqnarray}
which manifestly shows that the 2D problem reduces to computation of just flux--surfaces shape instead of usual
shape+labeling.

Notice that in the GSE equation, on the right--hand side, the quantity $H(x)$ has required both $g_2,g_3$
for its solution. However, both $g_2,g_3$ are also solutions of the 2D GSE equation. To converge these two
metric coefficients, iterations are necessary. Once convergence is reached, it is possible to check that the flux jump inside the plasma coincides with
the flux jump from current diffusion equation, i.e. $\Psi_{\rm b}-\Psi_0=\psi_{\rm b}-\psi_0$.

In cases of free--boundary computations, the plasma volume $V_{\rm b}$ can also change during iterations.
%
\subsection{The numerical scheme in details}
Astra produces input profiles to the equilibrium coupling schemes both $\psi(x)$ and $P(x)$. $x$ is here
the main Astra grid.
After entering the interface for the coupling the two profiles are preinterpolated on what will be used as
equilibrium grid $X$, equispaced from 0 to 1. As such, the iteration scheme sees entering profiles
$\psi(X),P(X)$. Note that usually $N_X < N_x$ to speed up equilibrium computations. 

Then the coupling scheme proceeds as follows (given for the general case of free--boundary equilibria):

0) $\psi(X),P(X)$ enter in the iteration process;

1ext) External iteration $j_{\rm ext}=0$: initial guesses for $g_2,g_3,\Phi_{\rm b},V_{\rm b}$ are set;

2ext) External iteration $j_{\rm ext}=j_{\rm ext}+1$: update $g_2^{j_{\rm ext}},g_3^{j_{\rm ext}},V_{\rm b}^{j_{\rm ext}}$;

1int) Internal iteration $i=0$: uses last computed $\Phi_{\rm b}$ as initial guess. 

2int) Internal iteration $i=i+1$: solve equation for $H(x)^{i,j_{\rm ext}}$ given $g_2^{j_{\rm ext}},g_3^{j_{\rm ext}},\Phi_{\rm b}^{i,j_{\rm ext}}$. Compute temporary volume $V_{\rm b}^{i}$. 

3int) If $V_{\rm b}^{i} \neq V_{\rm b}^{j_{\rm ext}}$ outside tolerance, adapt $\Phi_{\rm b}^{i,j_{\rm ext}}$ (done with bisection method), and go back to 2int). 

4int) $\Phi_{\rm b}^{i,j_{\rm ext}}$ has converged $\rightarrow$ compute $FF'$;

3ext) Call generic Grad--Shafranov solver with inputs $P'$ and $FF'$ as functions of $\psi_n=(\psi-\psi_0)/(\psi_{\rm b}-\psi_0)$. Returns outputs $g_2,g_3$ and new plasma volume $V_{\rm b}$;

4ext) Check relative error between old and new metric coefficients. If over tolerance, go back to 2ext);

1) Compute additional quantities requested by Astra;

2) Interpolate all outputs from equilibrium grid $X$ to transport grids $x,\tilde{x}$ accordingly.

After one such procedure, the program has performed $i$ iterations on 1D equation and $j_{\rm ext}$ 
iterations of 2D calculations. The output metric coefficients are self--consistent with given boundary or 
external coil currents and with current diffusion $\psi$.

Since all passages involve solution of parabolic equations on the same grid, no instability is present. If 
centered differencies are used, as in our case, spurious oscillations (pure imaginary Fourier modes) can be
present. However these are dealt with very easily by means of under--relaxation. In the present version,
old and new values are weighted 50$\%$ each to compute the updated value. This factor 0.5 is fixed for all 
computations and as such does not constitute a mean to force convergence, but simply to stabilize numerical
oscillations, retaining the 2nd order accuracy of the numerical solution of the equations. Our experience 
is that one needs to under--relax only $g_2$ and $V_{\rm b}$, which are the most sensisitve quantities to 
changes in plasma internal profiles.

Notice also that we kept this procedure very modular, such that any GS solver that uses P' and FF' as
inputs can be plugged into the iteration scheme.
%
\subsubsection{The numerical scheme: parameter namelist}
In the directory {\it exp/equ/(MACHIN)/} (see below for the meaning of this) there can be a file called {\it 
namelist.txt} that can be changed to try different numerical parameters. Here is the list with some example values 
(default values used by the code are optimized for basic calculations):
\\\\
0.002 $\:\:\:\:\:$     ! time from which to fix P' and FF' (ignores current diffusion of ASTRA)\\
0      $\:\:\:\:\:$    ! fix P' and FF' (ignores current diffusion of ASTRA)\\
5    $\:\:\:\:\:$      ! not used\\
0     $\:\:\:\:\:$      ! if value=1 writes an information file for testing equilbrium code standalone\\
0   $\:\:\:\:\:$         ! 0 - do iterations always, 1 - does iterations only at first FBE call, -1 - does iterations always except at first FBE call (recommended value is 0)\\
3   $\:\:\:\:\:$          ! numerical scheme for stabilization of numerical instabilities, 3 is recommended value\\
3    $\:\:\:\:\:$      ! hard-wired parameters, leave 3\\
0    $\:\:\:\:\:$      ! write files with diagnostic informations if value = 1\\
0    $\:\:\:\:\:$     ! leave 0\\
0.5   $\:\:\:\:\:$      ! under--relaxation factor\\
0.1    $\:\:\:\:\:$     ! under--relaxation factor for metric coefficients correction\\
0.2   $\:\:\:\:\:$    ! under--relaxation factor for metric coefficients derivatives\\
1.E+2  $\:\:\:\:\:$   ! not used, do not change\\
1.   $\:\:\:\:\:$      ! not used, do not change\\
10000   $\:\:\:\:\:$   ! maximum number of iterations allowed\\
1    $\:\:\:\:\:$    ! maximum number of iterations to be performed (if converges exits before of course). This depends on the time step. If time step is large, probably better to do some iterations\\
2    $\:\:\:\:\:$      ! interpolation routine: 1 is quadratic, 2 is cubic splines\\
4			  $\:\:\:\:\:$   			! interpolation method for flux--surfaces to rectangular grid: 4 is bi--quadratic\\
1.E-9    $\:\:\:\:\:$     ! tolerance on magnetic flux at boundary\\
1.E-5   $\:\:\:\:\:$      ! tolerance on metric coefficients variation\\
1.E-5    $\:\:\:\:\:$      ! tolerance on plasma volume variations\\
1.E-8 	  $\:\:\:\:\:$   		! tolerance on metric coefficients derivative variation (not used)\\					
0 	  $\:\:\:\:\:$   		! if 0, does not do iterations if refit currents. 1 to do iterations when refitting currents \\					
0 	  $\:\:\:\:\:$   		! if 0 does not overwrite coil.dat with refitted currents\\					

%
%\subsection{Another possible adiabatic compression scheme}
%In section 2.3 we have discussed how grid convective terms are treated in the present version when using
%normalized grid. There is also a third possibility, that is chosen by setting {\bf do\_adcmp=1}, or {\bf do\_adcmp=2}, %in the
%parameter file of the coupling scheme, i.e. in \\ {\it ./EQUIL/COUPLING\_SCHEME/parameters\_cpef2a.f90}%

%By setting {\bf do\_adcmp=1}, the code performs adiabatic compression dynamically during the internal 
%iterations of the 1D equation, such that $\psi(\Phi)$ remains an invariant. This is what is done on the 
%transport side a posteriori. However if done here, it assures consistency independently of the time 
%stepping.

%In addition, by setting {\bf do\_adcmp=2}, the same procedure is also applied to the pressure $P$, to keep $P(\Phi)$ constant (but neglecting $V'$ changes, which are anyway negligible in most cases).

%If instead {\bf do\_adcmp=0} is chosen, adiabatic compression is performed a posteriori in the transport 
%side. As said, that is still ok, but full consistency with the equilibrium requires solving it also in the 
%1D equation iterations.
%
\subsection{Choosing between free--boundary and prescribed boundary equilibria}
A new control parameter {\bf ITFBE} has been added, which can have values and meaning as explained in the following table 3.1.
\\
{\bf Table 3.1} Choice of equilibrium solution with {\bf ITFBE}\\[2ex]
\begin{tabular}{|l|l|l|}					\hline
\bf Name	&\bf Value & Description\\ \hline
\tt ITFBE 	& -1 & Equilibrium is solved in prescribed boundary mode\\
\tt  	& $>$ 0.0 &  Equilibrium is solved in free--boundary mode from TIME $>$ ITFBE\\ 
\hline
\end{tabular}
\bigskip\\
By default, {\bf ITFBE}=-1.

Another variable {\bf IFBEG} regulates the choice between rectangular grid with {\bf IFBEG=0} and adaptive grid with {\bf IFBEG=1}. Notice that this choice has also to be set in the equilibrium code as well otherwise some possible incomapitibility problems will arise.
%
\subsubsection{The SPIDER parameter list}
In the directory {\it exp/equ/(MACHIN)/} (see below for the meaning of this) there can be a file called {\it 
namelist$\_$spider.txt} that can be changed to try different numerical parameters. Here is the list with some example values 
(default values used by the code are optimized for basic calculations):
\\\\
0 $\:\:\:\:\:$     ! print info on screen from SPIDER\\
0      $\:\:\:\:\:$    ! choose k$\_$grid (0 for rectangular, 1 for adaptive)\\
1.0d-7    $\:\:\:\:\:$      ! tolerance on fixed boundary calculations\\
1.0d-6     $\:\:\:\:\:$      ! tolerance on circuit equations\\
1   $\:\:\:\:\:$         ! key$\_$plc: 1 to fix IPL into SPIDER, 0 to allow IPL from SPIDER to be used as actual current \\
4   $\:\:\:\:\:$          ! number of toroidal harmonics for TORIC\\
0    $\:\:\:\:\:$      ! if 1, save file for TORIC \\
0    $\:\:\:\:\:$      ! if 1, save file for TORBEAM \\
0    $\:\:\:\:\:$     ! if 1, save file for STRAHL \\
0   $\:\:\:\:\:$      ! if 1, save file for EQDSK \\
3    $\:\:\:\:\:$     ! number of fourier modes for STRAHL\\
0   $\:\:\:\:\:$    ! if > 0, save respective fort. file with some diagnostic data\\
%
\subsection{Free--boundary condition on FP equation}
In addition, the possibility of imposing an alternative boundary condition
to the current diffusion equation has been made available. The new control parameter
{\bf ITFBP} allows to choose this option. Its value are given in table 3.2:
\\
{\bf Table 3.2} Choice of boundary condition with {\bf ITFBP}\\[2ex]
\begin{tabular}{|l|l|l|}					\hline
\bf Name	&\bf Value & Description\\ \hline
\tt ITFBP 	& 0 & No effect (i.e. uses the ASTRA b.c., \\
\tt & & either prescribed $\ipl$ or prescribed $\psi_{\rm b}$)\\
\tt  	& -1 &  Boundary condition is $\displaystyle \psi_{\rm b}+\left[G \frac{\partial\psi}{\partial\rho}\right]_{\rm b}=\psi_{\rm ext}$\\ 
\hline
\end{tabular}
\bigskip\\
By setting {\bf ITFBP}=-1, the boundary condition $\displaystyle \psi_{\rm b}+\left[G \frac{\partial\psi}{\partial\rho}\right]_{\rm b}=\psi_{\rm ext}$ is used, where the left--hand--side is numerically discretized and the auxiliary values $G,\psi_{\rm ext}$ are provided by the equilibrium code from a free--boundary equilibrium computation. The exact definitions are: 
\begin{eqnarray}
G=\frac{1}{l_{\rm b}}\oint\oint \rm{Green}(\vec{r},\vec{r}') \; \left\vert\frac{\nabla{\rho}}{R}\right\vert(\vec{r}') \; dl \; dl' \\ \nonumber
\psi_{\rm ext}=\mu_0 \sum_j \; \frac{1}{l_{\rm b}}\oint  \rm{Green}(\vec{r}_j,\vec{r}') \; I_j \; dl'
\label{Gpsiextdefeq}
\end{eqnarray}
where "Green" is the Green function of the Grad--Shafranov Laplacian operator and $I_j$ are the external conductors currents, and $l_{\rm b}$ is the LCFS perimeter length.
The value of $\psib,\psi_{\rm ext},G$ in ASTRA are recorded in the scalar variables {\bf PSIFB}, {\bf PSIEXT}, {\bf PSPLEX}.
The default value of the b.c. control variable is {\bf ITFBP}=0.
%
\subsection{Choosing if to solve circuit equations and controller}
Additional control parameters {\bf ICIRCQ} and {\bf IPCTRL} have been introduced. The possible options are:
\\
{\bf Table 3.3} Choice of values for {\bf ICIRCQ} and {\bf IPCTRL}\\[2ex]
\begin{tabular}{|l|l|l|}					\hline
\bf Name	&\bf Value & Description\\ \hline
\tt ICIRCQ 	& 0 & Does not run circuit equations in equilibrium code \\
\tt  	& 1 &  Runs circuit equations in equilibrium code\\ 
\tt IPCTRL 	& 0 & Does not run controller, coil.dat untouched, no refit \\
\tt  	& 1 &  Runs controller, and refit coil--currents at first time step\\ 
\tt  	& -1 &  Runs controller, and solves FB--GSE at first time step, given coil currents\\ 
\tt  	& -2 &  No controller, do not refit coil--currents at first time step, \\
\tt  	&  &  CCOIL goes into coil.dat\\ 
\tt  	& -3 &  No controller, do not refit coil--currents at first time step,\\
\tt  	&  &   CCOIL goes into coil.dat\\ 
\tt  	& -4 &  No controller, refit coil--currents at first time step,\\
\tt  	&  &   coil.dat untouched\\ 
\tt  	& -5 &  No controller, refit coil--currents at first time step,\\
\tt  	&  &   CCOIL into coil.dat\\ 
\hline
\end{tabular}
\bigskip\\
Notice that if IPCTRL is non--integer, in particular if written as IPCTRL=X.1 (where X is a value as in the table), then
skips internal ASTRA iterations of NITREQ at first free--boundary call.

Notice also that if IPCTRL is such that CCOIL overwrites coil.dat, then one has to create, in the
{\it /exp/equ/(MACHINE)/} a directory called {\it req2d$\_$input} where a file called
{\it coil$\_$ref.dat} should be present. This file has the same structure as coil.dat, and will be used as a template 
to create the coil.dat using CCOIL values for the coil currents.
%
\subsection{Sign of toroidal field and toroidal current}
From Astra--7.0 on, it is explicitly stated that BTOR, IPL, and related parameters (CU, MU, etc.) have to be given with 
positive sign with respect to the reference coordinate system $(R,Z,\phi)$, or equivalently $(\psi,\theta,\phi)$. The 
sign as in the real experimental scenario to be simulated, has to be given separately with the two new control 
parameters {\bf SGNIP} and {\bf SGNBT}, as in table 3.3.
\\
{\bf Table 3.4} Sign of reference toroidal current and toroidal field\\[2ex]
\begin{tabular}{|l|l|l|}					\hline
\bf Name	&\bf Value & Description\\ \hline
\tt SGNIP 	& 1 & Current flows into positive $\phi$ direction\\
\tt  	& -1 & Current flows into negative $\phi$ direction\\ 
\tt SGNBT 	& 1 & Toroidal field is positive in $\phi$ direction\\
\tt  	& -1 &  Toroidal field is negative in $\phi$ direction\\ 
\hline
\end{tabular}
\bigskip\\
The signs will be taken into account in those parts of the code for which the direction of the field and current is relevant for the physical subsystem under study.
The default values are {\bf SGNIP}=1 and {\bf SGNBT}=1.
%
\subsection{PSIAX and PSIBO}
Two new output variables have been added to Astra. {\bf PSIAX} is the value of 
$\psi$ at $x=0$, exactly on the magnetic axis. Since the axis is not a grid point, 
{\bf PSIAX} is obtained with quadratic extrapolation from the first three grid 
points. {\bf PSIAX} is also used to compute the normalized poloidal grid properly.

The other variable {\bf PSIBO} is the value of $\psi$ at $x=1$.
%
\subsection{Machine dependency}
To allow trating multiple machines (e.g. AUG, ITER, JET, etc...), a new variable has been added to ASTRA.
The variable is called {\bf MACHIN} and it is a 4--character string that has to be assigned by putting the name of the machine in the first four characters of the experimental file list.
So if one wants to choose 'aug' the experimental file heading would read something like:\\\\
aug -AUG- \#11110 t=100 \\
 Param|Time  |  Value  |   Error \\
|-----|------|---------|-------- \\\\

For ITER it would look like \\\\
iter -ITER- \#11110 t=100 \\
 Param|Time  |  Value  |   Error \\
|-----|------|---------|-------- \\ \\

The variable {\bf MACHIN} is then assigned as 'aug ' or 'iter'.
The reason to do this is because, in the {\it exp/equ/} directory, the directory where equilibrium informations sit has the 
same name as trim(MACHIN) (i.e. without the eventual space in case of a 3--letters machine name). So for 'aug ' there 
will be a {\it exp/equ/aug/} directory which ASTRA and the equilibrium code will use to put information on the equilibrium.

These trick allows to avoid having to recompile the code every time one wants to switch between machines.







\clearpage





